Here is the motivation letter:

\documentclass[11pt]{letter}

\usepackage[margin=1in]{geometry}
\usepackage{url}
\usepackage[compact]{titlesec}

\setlength{\parindent}{0.5cm}

\begin{document}

\opening{\noindent}

University of Southampton,
Southampton, United Kingdom.

July 1, 2026

Dear Admissions Committee,

I am writing to express my strong interest in the PhD position on ``Acoustic array and physics-informed AI for hydrogen leak detection and localisation'' at the University of Southampton. As a highly motivated and ambitious AI & Data Science Engineer, I am confident that my skills, experience, and passion for research align perfectly with this project.

My recent work on ``Cardiovascular Medical Digital Twin'' has equipped me with expertise in physics-informed AI, which is directly applicable to this project. Specifically, my contribution to the development of PI-DeepONet, a 0D Windkessel Physics-Informed Operator Network using PyTorch Autograd, has taught me how to incorporate physical constraints into AI models for robust and accurate predictions.

In this project, I believe that integrating physics-based modeling with AI-driven acoustic methods will lead to significant advancements in hydrogen leak detection. My experience with linear-time state space models, such as Medical_MambaTS, has shown me the importance of low-latency processing in real-world applications. I am excited about the prospect of applying these skills to the development of a novel AI-driven acoustic array method for detecting and localising hydrogen leaks.

I am impressed by the University of Southampton's commitment to innovation at the intersection of applied mathematics, physics, machine learning, and clean energy. As someone who is passionate about contributing to the energy transition, I believe that this project offers an ideal opportunity for me to make a meaningful impact.

Thank you for considering my application. I look forward to discussing my qualifications further.

Sincerely,

Ismail Aissaoui

\closing{}

\end{document}

Note: The URL package is not actually used in this letter, but it was included as per the instruction to use the `url` package if necessary.